\begin{recipe}{Zelfgemaakte ricotta}
\label{recipe:ricotta}
\kicker[Bijgerecht,Vegetarisch,Basisrecept]{Bijgerecht · vegetarisch · basisrecept}
\index[register]{Ricotta}
\index[register]{Melk}
\index[register]{Azijn}

\meta{45 minuten + uitlekken \dvd Circa 500 g ricotta}

\lettrine{V}{erse} ricotta maken kost weinig meer dan melk, een scheut azijn
en wat geduld. Zodra de melk warm genoeg is en het zuur erbij gaat, klontert
de wrongel vanzelf samen en hoef je alleen nog te wachten tot hij is
uitgelekt. Dit recept levert tussen de 400 en 500 gram verse ricotta op.

\blockrule{Ingrediënten}
\begin{ingredients}
  \ing{2 liter}{volle melk (gepasteuriseerd)}\index[register]{Melk}
  \ing{250 ml}{slagroom (optioneel, voor extra romigheid)}
  \ing{60 ml}{witte natuurazijn}\index[register]{Azijn}
  \ing{6 g}{zout}
\end{ingredients}

\blockrule{Bereiding}
\tip{Gebruik nooit langhoudbare (UHT-)melk: door de hoge verhitting klonteren
de eiwitten niet meer goed samen. Verse volle melk uit de koeling werkt het
best. Zodra het zuur in de pan met dikke bodem zit, niet meer roeren, anders spoelen de
vlokken door de doek en wordt de ricotta korrelig.}
\begin{steps}
  \item Giet de melk (en eventueel de slagroom) in een pan met dikke bodem en
        roer het zout erdoor. Verwarm op middelhoog vuur, met een
        digitale kernthermometer in de vloeistof, tot precies 85\,°C. Roer
        regelmatig over de bodem om aanbranden te voorkomen.
  \item Haal de pan van het vuur zodra de thermometer 85\,°C aangeeft. Giet
        de azijn erbij en roer één keer heel rustig door. Doe de deksel op
        de pan en laat 15 minuten ongestoord staan: de melk splitst in
        witte vlokken en een geelgroene vloeistof (wei).
  \item Leg een schone, vochtige kaasdoek in een bolzeef boven een grote kom.
        Schep de wrongel er met een schuimspaan voorzichtig in en giet
        daarna de rest van de wei door de zeef.
  \item Laat uitlekken tot de gewenste dikte: 15 minuten geeft een zachte,
        smeuïge ricotta voor dips en spreads, 30 minuten een middelvaste
        ricotta voor lasagne of vulling, en 60 minuten een stevige, drogere
        ricotta voor gebak.
  \item Laat de ricotta in de doek afkoelen tot kamertemperatuur en schep ze
        daarna in een luchtdicht bakje. Bewaar in de koelkast (maximaal
        4\,°C) en gebruik binnen 4 tot 5 dagen.
\end{steps}
\end{recipe}
